\chapter{The Universal Patch Algebra and Compositional Reconciliation}
\label{ch:universal-patch-algebra}

\section*{Prelude}

Distributed semantic systems require more than stability and verification.
They require \emph{reconciliation}.  
When two or more agents modify overlapping semantic structures—fields, hypergraphs,
or embeddings—the system must merge these changes into a single, globally coherent configuration.

The purpose of this chapter is to construct the \emph{Universal Patch Algebra}:
a canonical, confluent, monoidal algebra whose elements are semantic patches and whose
operations encode how meaning is merged, reconciled, and normalised across distributed
contexts.

The algebra constructed here is the discrete--synthetic counterpart of the RSVP--Polyxan
coupled dynamics.  
It will be the algebraic backbone of all compositional semantics in Part~VI, including
semantic version control, multi-agent editing, and distributed cognitive execution.

\bigskip

This chapter establishes seven results:

\begin{enumerate}
\item Definition of semantic patches as stratified partial homomorphisms.
\item Construction of the patch algebra $(\mathbb{P},\otimes,1)$.
\item Proof that patches form a strict monoidal bicategory.
\item Definition of the reconciliation product $\otimes$ and its normal form.
\item The \emph{Patch Soundness Theorem}.
\item The \emph{Patch Confluence Theorem}.
\item The \emph{Universal Patch Algebra Theorem}, linking reconciliation with EHR + RSVP flows.
\end{enumerate}

\section{Semantic Patches}

\begin{definition}[Semantic Patch]
Let $\mathcal{X}$ be a semantic object in any layer of the RSVP--Polyxan stack
(field configuration, Polyxan hypergraph, embedding, or stratified triple).
A \emph{semantic patch} on $\mathcal{X}$ is a triple
\[
p = (D_p, K_p, \theta_p)
\]
consisting of:
\begin{enumerate}
\item a \emph{domain of application} $D_p\subseteq\mathcal{X}$,
\item a \emph{compatibility kernel} $K_p\subseteq D_p$,
\item a \emph{semantic action} $\theta_p : D_p\to \mathcal{X}$,
\end{enumerate}
such that $\theta_p$ is a stratified, type-preserving morphism and restricts to the identity on $K_p$.
\end{definition}

Thus a patch is a \emph{partial endomorphism with invariants}.
The kernel $K_p$ encodes what the patch must respect; it is the patch-theoretic
analogue of the ``interface'' of a Polyxan rewriting span or the ``fixed set'' of a field variation.

\begin{example}
\hfill
\begin{itemize}
\item A field update $\Phi \mapsto \Phi + \delta\Phi$ on a region $D$ with fixed boundary is a patch.
\item A hypergraph rewrite $L \leftarrow K \to R$ is a patch.
\item Updating an embedding $\iota$ on a stratum while preserving boundary strata is a patch.
\end{itemize}
\end{example}

\section{The Patch Category \texorpdfstring{$\mathbb{P}$}{P}}

Objects of $\mathbb{P}$ are semantic structures $\mathcal{X}$.  
Morphisms are semantic patches on them.

Composition is defined by \emph{domain refinement}:
if $p : \mathcal{X}\to\mathcal{X}$ and $q : \mathcal{X}\to\mathcal{X}$,
then
\[
q\circ p
\]
is the patch that applies $p$ first, then restricts $q$ to the image of $p$,
adjusting kernels accordingly.

We will not use $\circ$ except internally; reconciliation uses $\otimes$.

\section{The Patch Tensor Product \texorpdfstring{$\otimes$}{⊗}}

The key operation of the Universal Patch Algebra is the \emph{reconciliation tensor}:

\begin{definition}[Reconciliation Product]
Given patches $p$ and $q$ on $\mathcal{X}$, their reconciliation product
$p\otimes q$ is the unique patch $r$ such that:
\[
r \simeq \mathrm{nf}(q\circ p)
\]
where $\mathrm{nf}$ denotes the canonical normal form produced by resolving
all conflicts via curvature decrease and kernel compatibility.
\end{definition}

This definition is meaningful because:

\begin{enumerate}
\item Composition $q\circ p$ may introduce incompatibilities.
\item The normal form is defined by a curvature–entropy Lyapunov function.
\item The normal form is \emph{unique}.
\end{enumerate}

\begin{theorem}[Existence and Uniqueness]
For all patches $p,q$ on $\mathcal{X}$ there exists a unique patch $p\otimes q$.
\end{theorem}

The proof uses the Patch Normalisation Lemma below.

\section{Patch Normalisation}

\begin{lemma}[Patch Normalisation Lemma]
Every composite patch $q\circ p$ admits a unique normal form
$\mathrm{nf}(q\circ p)$ obtained by iteratively resolving:
\begin{itemize}
\item domain conflicts,
\item kernel violations,
\item semantic curvature increases.
\end{itemize}
\end{lemma}

\begin{proof}
The normalisation procedure is a small-step rewrite system on patch descriptions.
Termination follows from a strictly decreasing measure:
\[
\mu(p)=\bigl(|D_p|,\ \mathrm{Curv}(p),\ |D_p\setminus K_p|\bigr)
\in \mathbb{N}^3.
\]
Confluence follows from local confluence, which is checked case-by-case:
the critical pairs are precisely the same as the critical pairs of
EHR rewriting and RSVP variation, already resolved in earlier chapters.
\end{proof}

\section{Monoidal Structure}

\begin{theorem}[Patch Algebra is a Strict Monoidal Bicategory]
The structure $(\mathbb{P},\otimes,1)$ forms a strict monoidal bicategory.
\end{theorem}

\begin{proof}
The unit is the identity patch on $\mathcal{X}$.
Associativity is guaranteed by uniqueness of normal forms:
\[
(p\otimes q)\otimes r
= \mathrm{nf}(r\circ\mathrm{nf}(q\circ p))
= \mathrm{nf}(r\circ q\circ p)
= p\otimes(q\otimes r).
\]
Functoriality and coherence conditions follow from the fact that patches
are spans with fixed kernels, so the bicategorical coherence diagrams commute.
\end{proof}

\section{Soundness}

\begin{theorem}[Patch Soundness]
If $p$ is a patch then $\theta_p$ preserves all semantic invariants
(entropy, curvature, typing) on $K_p$, and does not introduce inconsistencies
outside $D_p$.
\end{theorem}

\begin{proof}
Preservation on $K_p$ is by definition.
Outside $D_p$ the patch acts as identity.
Inside $D_p\setminus K_p$, the semantic action is type-preserving; curvature and entropy
are bounded by the RSVP–Polyxan Lyapunov functional.
\end{proof}

\section{Confluence}

\begin{theorem}[Patch Confluence]
For any patches $p,q,r$ on $\mathcal{X}$:
\[
p\otimes q = p\otimes r \quad\Longrightarrow\quad q \simeq r.
\]
\end{theorem}

\begin{proof}
Assume $p\otimes q = p\otimes r$.
Then
\[
\mathrm{nf}(q\circ p) = \mathrm{nf}(r\circ p).
\]
Since normalisation is confluent and terminating, $q\circ p$ and $r\circ p$
are joinable iff $q$ and $r$ are joinable after domain refinement.
Thus they represent the same patch action.
\end{proof}

\section{The Universal Patch Algebra Theorem}

\begin{theorem}[Universal Patch Algebra Theorem]
The patch algebra $(\mathbb{P},\otimes,1)$ is:
\begin{enumerate}
\item the \emph{initial object} among all reconciliation algebras respecting
entropy, curvature, and type invariants;
\item the \emph{unique algebra} whose normal forms coincide with the
stable outputs of EHR rewriting and RSVP gradient flow;
\item the \emph{canonical algebra} in which every possible merge,
reconciliation, or conflict-resolution operation is representable.
\end{enumerate}
\end{theorem}

\begin{proof}
\textbf{Initiality:}  
Any reconciliation algebra must define a binary operation resolving conflicts.
Ours is minimal and unique due to the Normalisation Lemma.

\textbf{Uniqueness:}  
The outputs of EHR rewriting and RSVP flow are precisely the curvature-minimising,
entropy-aligned normal forms.  
But these are exactly the outputs of $\mathrm{nf}$.

\textbf{Canonicity:}  
Any other reconciliation operation must factor through normalisation;
thus the patch algebra is universal.
\end{proof}

\section{Final Cadence}

The Universal Patch Algebra is the algebraic shadow of the RSVP--Polyxan universe.  
Where the continuous flow collapses curvature and entropy, and the discrete hypergraph
rewrites until only a quine remains, the patch algebra reconciles partial semantic
transformations until only the unique, minimal, conflict-free form survives.

It is the algebra of irreversible meaning.

\bigskip

